
\include{./def/yf-formatting}
\include{./def/yf-def}

\usepackage{wasysym}
%\usepackage{times}

\title{}
\author{}
\date{}


%\def\A{\mathcal{A}}
%\def\B{\mathcal{B}}
\def\C{\mathcal{C}}
%\def\D{\mathcal{D}}
\def\E{\mathcal{E}}
\def\FF{\EuScript{F}}
\def\G{\mathcal{G}}
\def\GG{\EuScript{G}}
\def\H{\mathcal{H}}
\def\I{\mathcal{I}}
\def\J{\mathcal{J}}
\def\L{\mathcal{L}}
\def\Q{\mathcal{Q}}
\def\S{\mathcal{S}}
\def\T{\mathcal{T}}
\def\V{\mathcal{V}}
%\def\W{\EuScript{W}}
\def\W{\mathscr{W}}
\def\X{\mathcal{X}}
\def\Y{\mathcal{Y}}
\def\Z{\mathcal{Z}}
\def\att{\mit{diffatt}}
\def\agm{\mathit{AGM}}
\def\boldatt{\mathbf{att}}
\def\ext{\mathit{ext}}
\def\In{\mathrm{IN}}
\def\join{\mathit{join}}
\def\Out{\mathrm{OUT}}
\def\schema{\mathit{schema}}

\def\extraspacing{\vspace{4mm} \noindent}
\def\vgap{\vspace{4mm}}
\def\vslit{\vspace{2mm}}

\begin{document}

\boxminipg{\linewidth}{
    \begin{center}
        \vspace{3mm}
    {\Large
    (CE7456 Lecture Notes 12) The Geometry behind Natural Joins } \\[2mm]
    {\large Yufei Tao}
    \vspace{3mm}
    \end{center}
}

Today, we discuss a theorem that unveils an intrinsic mathematical structure underneath natural joins. This structure establishes a geometric, hierarchical view over the tuples in the join result, which enables simple algorithms for both computing the join and sampling from its result.

%\extraspacing {\bf Basic Definitions.}

\section{Basic Definitions}

Define $0^0 = 0$. For an integer $x \ge 1$, the notation $[x]$ denotes the set $\{1, 2, ..., x\}$. We use $\intdom$ to represent the set of integers.

\vgap

Let $\boldatt$ be a finite set, where each element is called an {\em attribute}. For a non-empty set $\X \subseteq \boldatt$ of attributes, a {\em tuple} over $\X$ is a function $\bm{u}: \X \rightarrow \intdom$; for each $X \in \X$, we refer to $\bm{u}(X)$ as the {\em value} of $\bm{u}$ on $X$. For any non-empty subset $\Y \subseteq \X$, we define $\bm{u}[\Y]$ --- called the {\em projection} of $\bm{u}$ on $\Y$ --- as the tuple $\bm{v}$ over $\Y$ satisfying $\bm{v}(Y) = \bm{u}(Y)$ for every attribute $Y \in \Y$. A {\em relation} $R$ is a set of tuples over the same set $Z$ of attributes; we refer to $Z$ as the {\em schema} of $R$ and represent it as $\schema(R)$.

\vgap

A {\em join} is defined as a set $\Q$ of relations. Let $\schema(\Q) = \bigcup_{R \in Q} \schema(R)$. The result of $\Q$ is a relation over $\schema(\Q)$ formalized as:
\myeqn{
    \join(\Q) &=& \{\textrm{tuple $\bm{u}$ over $\schema(\Q)$} \mid \forall R \in \Q: \bm{u}[\schema(R)] \in R\}. \label{eqn:join-rslt}
}
The {\em schema graph} of $\Q$ is a hypergraph $\H = (\V, \E)$ where
\myeqn{
    \V &=& \schema(\Q) \nn \\
    \E &=& \set{\schema(R) \mid R \in \Q}. \nn
}
Note that there is one-one correspondence between the relations in $\Q$ and the hyperedges in $\E$. For each hyperedge $e \in \E$, let $R_e$ represent the corresponding relation $R \in \Q$. Our discussion assumes that both $\V$ and $\E$ have constant sizes.

\vgap

A vector $\W$ over the space of $\E$ has a real-valued component $\W_e$ on each hyperedge $e \in \E$. The vector $\W$ is said to be a {\em fractional edge cover} of $\H$ if
\myitems{
    \item $0 \le \W_e \le 1$ for every $e \in \E$, and
    \item $\sum_{e \in \E : X \in e} \W_e \ge 1$ for every $X \in \V$.
}
Every fractional edge cover $\W$ of $\H$ is associated with an {\em AGM value} defined as
\myeqn{
    \mit{AGM}_\W &=& \prod_{e \in \E} |R_e|^{\W_e}. \label{eqn:agm}
}
As proved previously, it holds for any fractional edge cover $\W$ of $\H$ that
\myeqn{
    |\join(\Q)| \le \mit{AGM}_\W. \nn
}
The {\em AGM bound} of $\Q$ is the smallest AGM value of all fractional edge covers of $\H$.

\section{``Boxing'' Everything} \label{sec:split:box}

Set $d = |\V|$, and let
\myeqn{
    A_1, A_2, ..., A_d
    \label{eqn:topological}
}
be an (arbitrary) ordering of the attributes in $\V$. By treating each attribute as a {\em dimension}, we will often speak about the $d$-dimensional space $\intdom^d$, where $A_i$ ($i \in [d]$) corresponds to the $i$-th dimension. Note that each tuple in the join result $\join(\Q)$ can be regarded as a point in this space.

\vgap

We use the term {\em box} to refer to a rectangle $B$ in $\intdom^d$ of the form $[x_1, y_1] \times [x_2, y_2] \times ... \times [x_d, y_d]$. For any $i \in [d]$, let $B(A_i)$ represent the interval $[x_i, y_i]$, i.e., the projection of $B$ onto the $i$-th dimension. In particular, we say that $B(A_i)$ is a {\em singleton interval} if $x_i = y_i$. If $B(A_i)$ is a singleton interval on all $i \in [d]$, we call $B$ a  {\em singleton box}.

\vgap

For each relation $R \in \Q$, we define
\myeqn{
    R(B) = \{\text{tuple $\bm{u} \in R$} \mid
    \forall A_i \in \schema(R): \text{ $\bm{u}(A_i) \in B(A_i)$}\}.
    \label{eqn:R_B}
}
Intuitively, $R(B)$ is the subset of tuples in $R$ that are ``selected'' by $B$. Accordingly, define
\myeqn{
    \Q(B) = \set{R(B) \mid R \in \Q}.
    \label{eqn:Q_B}
}
Note that $\Q(B)$ is a join whose schema graph is $\H = (\V, \E)$, i.e., sharing the same schema graph as $\Q$.

\vgap

Take any fractional edge cover $\W$ of $\H$. The {\em AGM value of $\W$ on $B$} is defined as
\myeqn{
    \agm_{\W}(B) &=& \prod_{e \in \E} |R_e(B)|^{\W_e}.
    \label{eqn:box:agm}
}


\section{A Split Theorem} \label{sec:split}
%
% Given a non-singleton box $B$, we define its {\em split dimension} as the smallest $\sigma \in [d]$ such that $B(A_\sigma)$ is not a singleton interval. A set of boxes $B_1,$ $B_2,$ $..., B_s$ (where $s \ge 2$) is said to {\em refine} $B$ if
% \myitems{
%     \item $B_1, B_2, ..., B_s$ are mutually disjoint, and their union is $B$;
%     \item for every $i \in [d] \setm \set{\sigma}$, it holds that $B_1(A_i) = B_2(A_i) = ... B_s(A_i) = B(A_i)$.
% }

\begin{theorem} \textsc{(Simple AGM Split Theorem)} \label{thm:split-simple}
    Consider a join $\Q$ with schema graph $\H = (\V, \E)$ and an arbitrary fractional edge cover $\W$ of $\H$. Let $B$ be a non-singleton box; and let $B_1, B_2$ be any two disjoint boxes partitioning $B$ (i.e., $B_1 \cap B_2 = \emptyset$ and $B_1 \cup B_2 = B$). We must have
    \myeqn{
        \agm_W(B_1) + \agm_W(B_2) \le \agm_W(B). \nn
    }
\end{theorem}

\begin{proof}
   W.l.o.g., assume that $B_1$ and $B_2$ partition $B$ on dimension $A_1$, namely:
   \myitems{
        \item $B_1(A_1) \cap B_2(A_1) = \emptyset$ and $B_1(A_1) \cup B_2(A_1) = B(A_1)$;
        \item $B_1(A_i) = B_2(A_i)$ for all $i \in [2, d]$.
   }

   \vgap

   Let us review Friedgut's inequality. Fix arbitrary integers $p$ and $q$ at least 1. Take a set of non-negative real values $\{a_{i,j} \mid i \in [p], j \in [q]\}$. In addition, take another set of non-negative real values $\{b_i$ $\mid$ $i \in [q]\}$ satisfying $\sum_{i=1}^{q} b_i \geq 1$.  Friedgut's inequality \cite{f04} states that
    \myeqn{
        \sum_{i=1}^{p}	\prod_{j=1}^{q} a_{i, j}^{b_j}
        &\leq&
        \prod_{j=1}^{q} \left(\sum_{i=1}^{p} a_{i, j}\right)^{b_j}.
        \label{eqn:friedgut}
    }

    \vgap

    We can derive
    \myeqn{
        \agm_\W(B_1)
        &=&
        \prod_{e \in \E} |R_e(B_1)|^{\W_e} \nn \\
        &=&
        \Big(\prod_{e \in \E: A_1 \notin e} |R_e(B_1)|^{\W_e} \Big)
        \cdot
        \Big(\prod_{e \in \E: A_1 \in e} |R_e(B_1)|^{\W_e} \Big)
        \nn \\
        &=&
        \Big(\prod_{e \in \E: A_1 \notin e} |R_e(B)|^{\W_e} \Big)
        \cdot
        \Big(\prod_{e \in \E: A_1 \in e} |R_e(B_1)|^{\W_e} \Big)
        \label{eqn:split:1}
    }
    where the last equality used the fact that $R_e(B) = R_e(B_1)$ for every hyperedge $e \in \E$ that does not contain $A_1$. Similarly, a similar derivation yields
    \myeqn{
        \agm_\W(B_2)
        &=&
        \Big(\prod_{e \in \E: A_1 \notin e} |R_e(B)|^{\W_e} \Big)
        \cdot
        \Big(\prod_{e \in \E: A_1 \in e} |R_e(B_2)|^{\W_e} \Big).
        \label{eqn:split:2}
    }

    Since $\W$ is a fractional edge cover of $\H$, we have
    \myeqn{
        \sum_{e \in \E: A_1 \in e} \W_e \ge 1. \label{eqn:split:3}
    }
    For each hyperedge $e \in \E$, it holds that
    \myeqn{
        |R_e(B_1)| + |R_e(B_2)|
        &=&
        |R_e(B)|. \label{eqn:split:4}
    }
    To see why, note that each tuple of $R_e(B)$ appears in exactly one of $R_e(B_1)$ and $R_e(B_2)$ because $B$ is partitioned by $B_1$ and $B_2$.

    \vgap

    We therefore obtain
    \myeqn{
        && \agm_\W(B_1) + \agm_\W(B_2) \nn \\
        \explain{by \eqref{eqn:split:1} and \eqref{eqn:split:2}}
        &=&
        \Big(\prod_{e \in \E: A_1 \notin e} |R_e(B)|^{\W_e} \Big)
        \cdot
        \sum_{i=1}^2
        \prod_{e \in \E: A_1 \in e} |R_e(B_i)|^{\W_e} \nn \\
        \explain{by \eqref{eqn:friedgut} and \eqref{eqn:split:3}}
        &\le&
        \Big(\prod_{e \in \E: A_1 \notin e} |R_e(B)|^{\W_e} \Big)
        \cdot
        \prod_{e \in \E: A_1 \in e}
        \Big( \sum_{i=1}^2 |R_e(B_i)| \Big)^{\W_e} \nn \\
        \explain{by \eqref{eqn:split:4}}
        &\le&
        \Big(\prod_{e \in \E: A_1 \notin e} |R_e(B)|^{\W_e} \Big)
        \cdot
        \prod_{e \in \E: A_1 \in e}
        |R_e(B)|^{\W_e} \nn \\
        &=&
        \agm_\W(B) \nn
    }
    as claimed.
\end{proof}

\section{The Box-Tree Technique} \label{sec:boxtree}

Next, we use Theorem~\ref{thm:split-simple}  to design  algorithms. Let us fix
\myitems{
    \item a join $\Q$ with schema graph $\H = (\V, \E)$;
    \item an arbitrary attribute ordering $A_1, A_2, ..., A_d$ as defined in \eqref{eqn:topological} where $d = |\schema(\Q)|$;
    \item an arbitrary fractional edge cover $\W^*$ of $\H$ such that $\agm_{\W^*}$ equals the AGM bound of $\Q$.
}

Define
\myeqn{
    \In &=& \sum_{R \in \Q} |R|
}
which is the total number of tuples in the relations of $\Q$ (this is the {\em input size} of $\Q$).

\vgap

To focus our discussion on the mathematical core, we make four assumptions that allow us to abstract away implementation details:
\myitems{
    \item {{\bf A1:}} $\In$ is a power of 2.

    \item {{\bf A2:}} Attribute values lie in the so-called {\em rank space}. Specifically, for any relation $R \in \Q$, any attribute $A \in \schema(R)$, and any tuple $\bm{u} \in R$, we consider that $\bm{u}(A)$ is an integer in $[\In]$.

    \item {{\bf A3:}} Given any relation $R \in \Q$ and any box $B$, we can obtain $\agm_{\W^*}(B)$ --- see \eqref{eqn:box:agm} --- from an oracle for free.

    \item {{\bf A4:}} For any box $B$ with $\agm_{\W^*}(B) \le 1$, we can obtain the result of $\Q(B)$ in constant time.
}
%It is standard to satisfy these assumptions by preprocessing the relations of $\Q$ in $\tO(\In)$ time.

\vgap

We refer to the box $[\In]^d$ as the {\em join space}. By assumption {\bf A2}, every tuple of $\join(\Q)$ can be seen as a point in this space. The converse, however, does not hold: not every point in $[\In]^d$ necessarily corresponds to a tuple of $\join(\Q)$.


\subsection{The Box-Tree} \label{sec:boxtree:boxtree}

The {\em box tree} $\T$ is built by splitting the join space recursively. This is a binary tree whose nodes are each associated with a distinct box. Henceforth, we will denote a node using its associated box. For a node $B$, we define its {\em AGM value} as $\agm_{\W^*}(B)$; see \eqref{eqn:box:agm} for its calculation. The AGM value of a node is at least the sum of the AGM values of its children.

\vgap

The root of $\T$ is associated with the join space $[\In]^d$. In general, consider a node associated with box $B$. If $B$ has an AGM value at most $1$, we make $B$ a leaf of $\T$. Otherwise, $B$ must be a non-singleton box (the AGM value of every singleton box is at most 1), which we split \uline{\em evenly} along some arbitrary dimension into boxes $B_1$ and $B_2$. We then create two child nodes of $B$, associated with $B_1$ and $B_2$, respectively. This completes the definition of the box tree. It is worth mentioning that even splits are always possible due to assumption {\bf A1}.

\vgap

The tree $\T$ has four properties:
\myitems{
    \item {{\bf P1}}: Every internal node has an AGM value greater than $1$, while every leaf node has an AGM value at most $1$.
    \item {{\bf P2}}: The leaf boxes are mutually disjoint and their union is $[\In]^d$.
    \item {{\bf P3}}: $\T$ has height $O(\log \In)$.
    \item {{\bf P4}}: $\T$ has $O(\agm_{\W^*} \cdot \log \In)$ nodes.

}
Properties {\bf P1} and {\bf P2} follow immediately from the definition of $\T$. To understand {\bf P3}, notice that the box of a child node covers half as many points in the join space as that of its parent node. Since the join space has $\In^d$ points (assumption {\bf A2}), it can be split at most $d \log \In = O(\log \In)$ times before reducing to a singleton box. Hence, the height of $\T$ is $O(\log \In)$.

\vgap

To prove {\bf P4}, imagine removing all the leaf nodes of $\T$, and let $\T'$ denote the resulting tree (note that $\T'$ need not be binary). Every leaf of $\T'$ is an internal node in $\T$ and hence has an AGM value greater than $1$. The sum of the AGM values of all leaves of $\T'$ cannot exceed the AGM value of the root, which is precisely $\agm_{\W^*}$. Hence, $\T'$ has less than $\agm_{W^*}$ leaves, and thus at most $O(\agm_{W^*} \cdot \log \In)$ nodes in total. This proves {\bf P4}, since each node of $\T'$ can parent at most 2 leaves of $\T$.

\subsection{Join Reporting} \label{sec:boxtree:report}

Property {\bf P2} of the box tree $\T$ implies
\myeqn{
    \bigcup_{\text{leaf } B} \join(\Q(B)) &=& \join(\Q). \nn
}
where $\Q(B)$ is defined in \eqref{eqn:Q_B}. To compute $\join(\Q)$, generate the entire $\T$, which (by assumption {\bf A3} and property {\bf P4}) requires $O(\agm_{\W^*} \log \In)$  time. By property {\bf P1}, each leaf $B$ of $\T$ has an AGM value at most 1; hence, we can compute $\join(\Q(B))$ in $O(1)$ time (assumption {\bf A4}). This yields a join algorithm with running time $O(\agm_{\W^*} \cdot \log \In)$.

\subsection{Join Sampling} \label{sec:boxtree:sample}

The box-tree technique also gives a simple algorithm for {\em join sampling}, where the goal is to draw a tuple from $\join(\Q)$ uniformly at random. For this purpose, we do not materialize $\T$, but treat join sampling as a {\em tree sampling} problem discussed earlier in the course. The details are left to you. We discuss here only the special case where $\Out = 0$, because  tree sampling requires infinite repeats in this scenario. To remedy the issue, we stop tree sampling after $\agm_{\W^*}$ repeats and revert to the algorithm in Section~\ref{sec:boxtree:report} for computing $\join(\Q)$ in $O(\agm_{\W^*} \cdot \log \In)$ time, which either confirms $\Out = 0$ or allows us to obtain a random tuple of $\join(\Q)$ at no extra cost. The overall sampling time is
\myeqn{
    O\Big(
        \fr{\agm_{\W^*}}{\max\{1, \Out\}} \cdot \log \In
    \Big)
    \nn
    %\label{eqn:boxtree:sample-time}
}
in expectation where $\Out = |\join(\Q)|$.


\section{Remark}

Theorem~\ref{thm:split-simple} is a special case of a result proved in \cite{tw25}. The box-tree technique was developed in  \cite{dlt23}.

\bibliographystyle{plainurl}% the mandatory bibstyle
\bibliography{ref}



\end{document}
